\chapter{Methodology}
\label{ch:methodology}

This chapter describes the design and implementation of the proposed
Graph Convolutional Anatomically-constrained Dual-supervised Architecture
(GCADA) for clinical rehabilitation pose assessment. As discussed in
Chapter~\ref{ch:literature}, prior whole-body 3D pose estimation methods
treat the body, face, and hands as independent regions and supervise
exclusively on positional error in millimetres, leaving the clinically
relevant degree-valued range-of-motion (ROM) signal entirely untouched.
GCADA closes this gap by coupling the multi-resolution HR-GCN backbone
of Zhang et al.~\cite{zhang2025hrgcn} to two additional learning
mechanisms that are introduced for the first time in this thesis: an
\textbf{Anatomical Constraint Layer} that penalises predictions falling
outside clinically valid ROM windows, and a \textbf{Clinical Angle
Supervision Loss} that supervises a lightweight regression head directly
on Vicon-measured joint angles in degrees. End-to-end, the system maps
a monocular RGB stream through seven stages: (i) RTMPose-W 2D keypoint
detection~\cite{jiang2023rtmpose}, (ii) HR-GCN 3D lifting, (iii) per-frame
ROM regression by the proposed angle head, (iv) hard anatomical clamping
at inference, (v) One-Euro temporal smoothing~\cite{casiez2012euro},
(vi) geometric ROM verification, and (vii) clinical reporting.

\section{System Overview}
\label{sec:system-overview}

The proposed pipeline operates in two distinct modes that share a
common forward graph. During \emph{training}, the network ingests
pre-extracted 2D keypoints together with paired Vicon-derived 3D
positions and ROM angles from the UI-PRMD dataset, optimising a
composite loss that jointly supervises millimetre position error,
degree angle error and an anatomical-validity penalty. During
\emph{inference}, raw video frames are passed through the RTMPose-W
detector to obtain 133 COCO-WholeBody 2D keypoints, the HR-GCN
backbone lifts these to 3D, the ClinicalAngleHead regresses 12 ROM
angles, the angles are hard-clamped to physiologically valid ranges,
and the 3D joint trajectories are smoothed by per-coordinate One-Euro
filters before being passed to the geometric ROM verifier and the
clinician-facing report. A subject-level split is enforced so that no
patient appearing during training is ever observed at evaluation,
ensuring that the reported ROM error reflects generalisation to unseen
individuals rather than memorisation of training subjects.

\section{Dataset: UI-PRMD}
\label{sec:dataset}

The University of Idaho Physical Rehabilitation Movements Dataset
(UI-PRMD) of Vakanski et al.~\cite{vakanski2018uiprmd} is used for all
fine-tuning and evaluation. UI-PRMD contains optical-marker recordings
captured by a 39-marker Vicon system at 60~Hz, providing ground-truth
joint positions in millimetres and joint angles in degrees. Ten
exercises (m01--m10) are performed by ten healthy subjects (s01--s10);
the recording for subject~3, exercise~3 is missing. The subject-level
split adopted by this work, train: s01--s08 and test: s09--s10, follows
a leave-subjects-out protocol: it eliminates within-subject leakage
and is the standard generalisation regime for clinical biomechanics.

The raw \texttt{m\{ee\}\_s\{ss\}\_positions.txt} files are converted
into NumPy archives by the preprocessor in
\texttt{utils/prepare\_data\_uiprmd.py}. Each Vicon frame yields three
synchronised tensors: a 133-joint COCO-WholeBody 2D pose
$\mathbf{P}^{2D}\!\in\!\mathbb{R}^{133\times 2}$ obtained by an
orthographic frontal projection $(X,-Z)$ of the marker positions;
a 133-joint 3D pose
$\mathbf{P}^{3D}\!\in\!\mathbb{R}^{133\times 3}$ in which the 23 body
joints are populated from the Vicon-to-COCO mapping documented in
the preprocessor and the 68 face and 42 hand joints are zero-padded;
and a vector $\boldsymbol{\theta}^{*}\!\in\!\mathbb{R}^{12}$ of clinical
ROM angles in degrees, computed geometrically from the marker positions
(Section~\ref{sec:rom-extraction}). Both 2D and 3D poses are
hip-centred at the midpoint of the LASI/RASI markers; 2D coordinates
are divided by a fixed physical scale of $1\,000$~mm so that their
empirical standard deviation matches that of the H3WB pre-training
distribution, and 3D coordinates are converted to metres. After
preprocessing the splits contain $204\,194$ training frames and
$49\,475$ test frames.

\section{2D Pose Detection via RTMPose-W}
\label{sec:rtmpose}

The first stage of the inference pipeline detects 2D keypoints on each
incoming RGB frame. RTMPose-W~\cite{jiang2023rtmpose} is selected for
its real-time performance and its native support for the COCO-WholeBody
keypoint definition that the HR-GCN backbone expects. It produces 133
keypoints per person comprising 17 body joints, 68 facial landmarks,
42 hand joints (21 per hand) and 6 foot keypoints. The detector
output is centred on the hip midpoint and divided by the same physical
scale used during training, yielding 2D coordinates in the same
normalised range $[-1,1]$ as the synthetic projections that drive
training; this matched preprocessing avoids a domain gap between the
training-time 2D distribution and the deployment-time detector output.

\section{HR-GCN Backbone for 3D Lifting}
\label{sec:hrgcn}

The 2D-to-3D lifting backbone is the multi-branch HR-GCN of Zhang et
al.~\cite{zhang2025hrgcn} (\texttt{models/graph\_hrnet\_multi\_branch.py},
configuration \texttt{w32\_adam\_lr1e-3.yaml}, model variant 1).
The network combines a High-Resolution Module (HRM) that maintains
parallel multi-resolution graph branches with a Fine-Grained Keypoint
Module (FGKPM) that processes body, face, left hand and right hand on
dedicated heads, and is initialised from the public checkpoint
pre-trained on the H3WB dataset~\cite{Zhu2023H3WB} (reported MPJPE of
57.4~mm on H3WB). Given an input
$\mathbf{P}^{2D}\!\in\!\mathbb{R}^{B\times 133\times 2}$ the backbone
emits four part-specific 3D tensors that are concatenated to recover a
full 3D whole-body pose
$\hat{\mathbf{P}}^{3D}\!\in\!\mathbb{R}^{B\times 133\times 3}$.

The graph convolutional operator used throughout this work is the
\texttt{dc\_preagg} variant of Decoupled Pre-Aggregation Graph
Convolution. For an input feature $\mathbf{H}\!\in\!\mathbb{R}^{J\times C}$
with adjacency $\mathbf{A}\!\in\!\mathbb{R}^{J\times J}$ and identity
$\mathbf{I}$, the layer applies separate joint-wise weight tensors
$\mathbf{W}_{0}, \mathbf{W}_{1}\!\in\!\mathbb{R}^{J\times C\times C'}$
to the self-loop and the neighbour aggregation respectively,
\begin{equation}
\mathbf{H}' = (\mathbf{A}\odot\mathbf{I})\,(\mathbf{H}\,\mathbf{W}_{0})
            + (\mathbf{A}\odot(\mathbf{1}-\mathbf{I}))\,(\mathbf{H}\,\mathbf{W}_{1}) + \mathbf{b},
\label{eq:dcpreagg}
\end{equation}
where $\odot$ denotes the Hadamard product. Equation~\ref{eq:dcpreagg}
generalises the standard symmetric propagation rule of
Kipf and Welling~\cite{kipf2017semi} by decoupling the parameters that
govern a node's self-update from those that govern its aggregation
from neighbours, which is essential for skeletons whose joints play
heterogeneous functional roles. During fine-tuning on UI-PRMD, the
backbone parameters are updated at half the head learning rate
(\texttt{backbone\_lr\_factor}~$=0.5$) and are temporarily frozen for
the first three epochs while the freshly initialised angle head warms
up.

\section{Novel Contribution 1: Proposed Anatomical Constraint Layer}
\label{sec:n1-anatomical}

\subsection{Motivation}
A model that minimises only positional error has no incentive to
respect physiological joint limits, and in practice may emit elbow
hyperextensions, knee bend reversals or trunk torsions that are
physically impossible to produce in a healthy human body. While prior
work such as Dabral et al.~\cite{Dabral2018ECCV} introduced a generic
illegal-angle penalty for action-recognition oriented pose, the
present work is the first to embed a \emph{rehabilitation-specific}
ROM constraint into a whole-body pose estimator and to enforce it both
as a soft training penalty and as a hard inference clamp.

\subsection{Joint ROM Limits}
The constraint vocabulary spans the twelve clinically actionable ROM
joints used throughout UI-PRMD. The minimum and maximum admissible
angles in degrees, encoded in the dictionary
\texttt{REHAB\_JOINT\_LIMITS}, are listed in
Table~\ref{tab:rom-limits}.

\begin{table}[h]
\centering
\caption{Clinical ROM limits enforced by the proposed Anatomical
Constraint Layer. Ankle convention: $90^{\circ}$ is anatomical neutral;
values above $90^{\circ}$ denote dorsiflexion and below $90^{\circ}$
plantarflexion.}
\label{tab:rom-limits}
\begin{tabular}{lcc}
\hline
Joint & Min ($^{\circ}$) & Max ($^{\circ}$) \\
\hline
Cervical pitch          & $-60$  & $60$  \\
Trunk flexion           & $-90$  & $90$  \\
Left  shoulder flexion  & $-30$  & $180$ \\
Right shoulder flexion  & $-30$  & $180$ \\
Left  shoulder abduction& $0$    & $180$ \\
Right shoulder abduction& $0$    & $180$ \\
Left  hip               & $-30$  & $120$ \\
Right hip               & $-30$  & $120$ \\
Left  knee              & $0$    & $150$ \\
Right knee              & $0$    & $150$ \\
Left  ankle             & $50$   & $130$ \\
Right ankle             & $50$   & $130$ \\
\hline
\end{tabular}
\end{table}

\subsection{Training: Soft Quadratic Penalty}
Let $\hat{\theta}_{j}$ denote the predicted ROM angle for joint
$j\!\in\!\{1,\dots,J\}$ with $J\!=\!12$, and let $(\theta^{\min}_{j},
\theta^{\max}_{j})$ be its admissible interval. The proposed soft
penalty is the mean squared excursion outside that interval,
\begin{equation}
\mathcal{L}_{\mathrm{constraint}} \;=\;
\frac{1}{J}\sum_{j=1}^{J}
\Big[\,\max\!\big(0,\;\theta^{\min}_{j}-\hat{\theta}_{j}\big)
+ \max\!\big(0,\;\hat{\theta}_{j}-\theta^{\max}_{j}\big)\,\Big]^{2}.
\label{eq:constraint}
\end{equation}
The loss is identically zero whenever the prediction lies inside the
valid ROM window and grows quadratically with the magnitude of the
violation, providing a smooth gradient that pushes physiologically
implausible predictions back into the admissible region.

\subsection{Inference: Hard Clamp}
At deployment, predictions are passed through a hard element-wise
clamp,
\begin{equation}
\hat{\theta}^{\mathrm{out}}_{j} = \mathrm{clip}\!\big(\hat{\theta}_{j},\,
\theta^{\min}_{j},\,\theta^{\max}_{j}\big),
\label{eq:clamp}
\end{equation}
implemented in \texttt{clamp\_angles\_to\_valid\_range}. This
guarantees that every angle reported to a clinician is physiologically
valid by construction, irrespective of any residual numerical error in
the regression head.

\section{Novel Contribution 2: Proposed Clinical Angle Supervision Loss}
\label{sec:n2-clinical-loss}

\subsection{ClinicalAngleHead Architecture}
A lightweight multi-layer perceptron, denoted \texttt{ClinicalAngleHead},
maps the 23 body joints emitted by the HR-GCN backbone to the 12
clinical ROM angles. Given the body sub-tensor
$\mathbf{B}\!\in\!\mathbb{R}^{B\times 23\times 3}$, the head flattens
each sample to 69 features and applies the sequence
\begin{equation}
\hat{\boldsymbol{\theta}} \;=\;
\mathbf{W}_{3}\,\sigma\!\big(\mathbf{W}_{2}\,
\mathrm{Drop}_{0.1}\!\big(\sigma(\mathbf{W}_{1}\mathbf{b})\big)\big),
\label{eq:anglehead}
\end{equation}
with $\sigma\!=\!\mathrm{ReLU}$, hidden widths $69\!\rightarrow\!128
\!\rightarrow\!64\!\rightarrow\!12$ and a dropout rate of $0.1$ between
the first two hidden layers. The complete head contains $17\,996$
trainable parameters, which is negligible compared to the HR-GCN
backbone, but is sufficient to recover the non-linear mapping from
joint coordinates to ROM angles. The output ordering follows
\texttt{JOINT\_LIMIT\_TENSOR\_ORDER} so that the constraint penalty
of Equation~\ref{eq:constraint} can be applied directly without
re-indexing.

\subsection{Combined Loss Function}
The total objective is a weighted sum of the standard whole-body
position loss, the proposed clinical angle loss and the proposed
anatomical constraint penalty,
\begin{equation}
\mathcal{L}_{\mathrm{total}} \;=\;
\mathcal{L}_{\mathrm{pos}} \;+\;
\lambda_{\mathrm{angle}}\,\mathcal{L}_{\mathrm{angle}} \;+\;
\lambda_{\mathrm{constraint}}\,\mathcal{L}_{\mathrm{constraint}},
\label{eq:total-loss}
\end{equation}
where the position term reproduces the per-part Mean-Squared-Error /
$L_{1}$ blend of the original HR-GCN,
\begin{equation}
\mathcal{L}_{\mathrm{pos}} = \!\!\sum_{p\in\{\mathrm{body},\mathrm{face},\mathrm{lhand},\mathrm{rhand}\}}\!\!
w_{p}\,\Big[(1-\mu)\,\mathrm{MSE}(\hat{\mathbf{P}}_{p},\mathbf{P}_{p})
+ \mu\,L_{1}(\hat{\mathbf{P}}_{p},\mathbf{P}_{p})\Big],
\label{eq:pos-loss}
\end{equation}
with $\mu\!=\!0.5$ and part weights $w_{\mathrm{body}}\!=\!1.0$,
$w_{\mathrm{face}}\!=\!0.5$, $w_{\mathrm{hand}}\!=\!0.5$. The clinical
angle term is the mean absolute error between predicted and
Vicon-measured ROM angles in degrees,
\begin{equation}
\mathcal{L}_{\mathrm{angle}} \;=\;
\frac{1}{J}\sum_{j=1}^{J} \big|\hat{\theta}_{j} - \theta^{*}_{j}\big|,
\label{eq:angle-loss}
\end{equation}
and $\mathcal{L}_{\mathrm{constraint}}$ is the soft penalty of
Equation~\ref{eq:constraint}. The full GCADA configuration sets
$\lambda_{\mathrm{angle}}\!=\!0.1$ and
$\lambda_{\mathrm{constraint}}\!=\!0.05$; the baseline ablation that
isolates the contribution of the angle supervision sets both weights
to zero.

\subsection{Why this is novel}
No previously published whole-body pose estimator supervises ROM
angles directly in degree units. Physio2.2M~\cite{Rode2025SciReports}
provides large-scale physiotherapy data but its accompanying baselines
optimise only positional error. HGcnMLP~\cite{Hu2024_FrontBioengBiotechnol_HGcnMLP}
introduces an MLP refinement on top of a GCN but again supervises in
millimetres and computes ROM only as a downstream evaluation metric.
The combination of (i) end-to-end angle supervision in degrees,
(ii) a soft anatomical constraint penalty and (iii) a hard inference
clamp, all operating on a whole-body 3D pose estimator, is original
to this thesis.

\section{Temporal Smoothing via the One-Euro Filter}
\label{sec:one-euro}

Per-frame predictions are temporally noisy when fed by a real-time 2D
detector. To suppress jitter without introducing perceptible lag,
GCADA applies the One-Euro filter of Casiez et al.~\cite{casiez2012euro}
independently to every joint coordinate. The filter adapts its cut-off
frequency to the instantaneous signal speed: at low velocities a low
cut-off heavily smooths the signal, while at high velocities the
cut-off rises so that fast movements are not lagged. For sample $x_{t}$
at time $t$ with sampling rate $f_{s}$, the filtered velocity estimate
$\hat{\dot{x}}_{t}$ and signal $\hat{x}_{t}$ satisfy
\begin{equation}
\hat{\dot{x}}_{t} = \alpha(f_{c}^{d})\,(x_{t}-\hat{x}_{t-1})f_{s}
+(1-\alpha(f_{c}^{d}))\,\hat{\dot{x}}_{t-1},
\quad
f_{c} = f_{c}^{\min} + \beta\,|\hat{\dot{x}}_{t}|,
\label{eq:euro-vel}
\end{equation}
\begin{equation}
\hat{x}_{t} = \alpha(f_{c})\,x_{t} + (1-\alpha(f_{c}))\,\hat{x}_{t-1},
\quad
\alpha(f) = \frac{1}{1+\tau(f) f_{s}},
\quad
\tau(f) = \frac{1}{2\pi f}.
\label{eq:euro-pos}
\end{equation}
The implementation in \texttt{common/one\_euro\_filter.py} uses
$f_{c}^{\min}\!=\!1.0$~Hz, $\beta\!=\!0.007$, $f_{c}^{d}\!=\!1.0$~Hz
and the deployment frame rate $f_{s}\!=\!30$~Hz. The wrapper class
\texttt{SkeletonFilter} instantiates one independent filter per
coordinate, yielding $133\!\times\!3=399$ stateful filters that
together smooth the full whole-body trajectory before geometric
verification.

\section{Clinical ROM Angle Extraction}
\label{sec:rom-extraction}

The geometric ROM extractor in
\texttt{utils/prepare\_data\_uiprmd.py} is used both to generate the
training labels $\boldsymbol{\theta}^{*}$ from raw Vicon data and to
verify the ClinicalAngleHead predictions at inference time. All angles
are computed in vectorised form using the basic identity
\begin{equation}
\angle(\mathbf{u},\mathbf{v}) = \arccos\!\left(\frac{\mathbf{u}\cdot\mathbf{v}}
{\lVert\mathbf{u}\rVert\,\lVert\mathbf{v}\rVert}\right).
\label{eq:angle}
\end{equation}
With $\hat{\mathbf{z}}$ the Z-up vertical axis and the spine vector
defined as $\mathbf{s}=\mathbf{m}_{\mathrm{sho}}-\mathbf{m}_{\mathrm{hip}}$,
the twelve clinical conventions implemented are
\begin{align}
\theta_{\mathrm{cerv}}     &= \angle(\mathbf{p}_{\mathrm{head}}-\mathbf{p}_{\mathrm{C7}},\,\hat{\mathbf{z}}), \\
\theta_{\mathrm{trunk}}    &= \angle(\mathbf{s},\,\hat{\mathbf{z}}), \\
\theta^{L}_{\mathrm{shoF}} &= 180^{\circ} - \angle(\mathbf{s},\,\mathbf{p}_{\mathrm{Lelb}}-\mathbf{p}_{\mathrm{Lsho}}), \\
\theta^{L}_{\mathrm{shoA}} &= 180^{\circ} - \angle(\hat{\mathbf{z}},\,\Pi_{\mathrm{frontal}}(\mathbf{p}_{\mathrm{Lelb}}-\mathbf{p}_{\mathrm{Lsho}})), \\
\theta^{L}_{\mathrm{hip}}  &= 180^{\circ} - \angle(\hat{\mathbf{z}},\,\mathbf{p}_{\mathrm{Lkne}}-\mathbf{p}_{\mathrm{Lhip}}), \\
\theta^{L}_{\mathrm{knee}} &= 180^{\circ} - \angle(\mathbf{p}_{\mathrm{Lhip}}-\mathbf{p}_{\mathrm{Lkne}},\,\mathbf{p}_{\mathrm{Lank}}-\mathbf{p}_{\mathrm{Lkne}}), \\
\theta^{L}_{\mathrm{ankle}}&= \angle(\mathbf{p}_{\mathrm{Lank}}-\mathbf{p}_{\mathrm{Lkne}},\,\mathbf{p}_{\mathrm{Ltoe}}-\mathbf{p}_{\mathrm{Lhee}}),
\label{eq:rom-defs}
\end{align}
and analogously for the right side. The shoulder abduction projects
the upper-arm vector onto the frontal plane via
$\Pi_{\mathrm{frontal}}(\mathbf{u})=\mathbf{u}-(\mathbf{u}\!\cdot\!
\mathbf{n})\mathbf{n}$, where $\mathbf{n}$ is the unit normal of the
frontal plane obtained as $\mathbf{n}=(\mathbf{p}_{\mathrm{Rsho}}
-\mathbf{p}_{\mathrm{Lsho}})\times\hat{\mathbf{z}}$. Knee and hip
flexion are subtracted from $180^{\circ}$ so that $0^{\circ}$
corresponds to full extension; the ankle uses the heel-to-toe vector
as the foot's long axis so that anatomical neutral evaluates to
$90^{\circ}$.

\section{Training Configuration}
\label{sec:training}

Training is launched by the script \texttt{train\_rehab.py} and uses
the H3WB-pretrained checkpoint as the backbone initialisation.
Hyper-parameters are summarised in Table~\ref{tab:training-hp}.

\begin{table}[h]
\centering
\caption{Training configuration used for the proposed GCADA pipeline.}
\label{tab:training-hp}
\begin{tabular}{lc}
\hline
Hyper-parameter & Value \\
\hline
Optimiser                            & Adam \\
Base learning rate (angle head)      & $5\!\times\!10^{-4}$ \\
Backbone learning-rate factor        & $0.5$ \\
Backbone learning rate (effective)   & $2.5\!\times\!10^{-4}$ \\
Mini-batch size                      & $64$ \\
Epochs                               & $50$ \\
LR scheduler                         & ReduceLROnPlateau ($\times 0.5$, patience $8$, $\min 10^{-6}$) \\
Backbone warm-up (frozen) epochs     & $3$ \\
Gradient clipping (max-norm)         & $1.0$ \\
$\lambda_{\mathrm{angle}}$ (full GCADA)     & $0.1$ \\
$\lambda_{\mathrm{constraint}}$ (full GCADA)& $0.05$ \\
$\lambda_{\mathrm{angle}}$ / $\lambda_{\mathrm{constraint}}$ (baseline) & $0.0$ / $0.0$ \\
Position-loss MSE/$L_{1}$ blend $\mu$       & $0.5$ \\
GCN variant                          & \texttt{dc\_preagg} \\
Backbone model                       & HR-GCN Model 1 (multi-branch) \\
Mirror augmentation probability      & $0.5$ (training) \\
GPU                                  & NVIDIA A10G (24~GB) \\
\hline
\end{tabular}
\end{table}

The optimiser uses parameter groups so that the H3WB-pretrained
backbone receives a learning rate two times lower than the freshly
initialised angle head, balancing rapid adaptation of the new head
with conservative refinement of the well-trained backbone. During
the first three epochs the backbone is frozen entirely so that the
random angle head can stabilise without corrupting the lifting
representation. An optional progressive weighting schedule, used
during early experimentation, ramps both $\lambda_{\mathrm{angle}}$
and $\lambda_{\mathrm{constraint}}$ from $0.1\times$ to $1.0\times$
of their nominal values across the training run, which proved useful
for the larger H3WB-to-Vicon domain shift. A horizontal mirror
augmentation is applied with probability $0.5$ during training; it
flips the X coordinate of the 2D input, swaps the left/right body
joints in the 3D target and exchanges the corresponding left/right
ROM angles, doubling the effective training set without altering the
underlying biomechanics.

Partial weight loading is used to preserve compatibility with earlier
six-angle checkpoints. When a checkpoint produced by an earlier
six-output ClinicalAngleHead is provided, the backbone is loaded in
full and the hidden layers of the angle head are copied verbatim;
the first six rows of the output layer are inherited and the new
rows corresponding to the six newly added ROM joints are left at
their random initialisation. This lets the project evolve its
clinical vocabulary without discarding earlier training progress.

The best checkpoint is selected by the lowest mean ROM mean-absolute
error on the held-out s09--s10 evaluation split, computed at the end
of every epoch. The same evaluation routine also reports body, face
and hand MPJPE in millimetres so that any improvement in clinical
angle accuracy is contextualised against the whole-body positional
quality preserved by the original HR-GCN objective.
